% Evans 6.6: English printed pp.343--348 / PDF pp.360--365.
\Needspace{65mm}
\section{Problems and Worked Solutions}
\label{sec:6-6}
These are concise restatements of the 17 source problems, with editorial
solutions. Unless a problem says otherwise, $U$ is bounded with smooth
boundary, coefficients are smooth and bounded, and ellipticity is uniform.
Connectedness is specified where constants or positivity are involved.
The same assumptions apply to each weak formulation below.

\Needspace{30mm}
\paragraph{1. Removing a potential by a positive solution.}
Relate $-\Delta u+cu=0$ to $-\operatorname{div}(aDv)=0$.
\begin{proof}[Solution]
If $w>0$ solves the first equation, set $u=wv$. The product rule gives
\[
 -\Delta(wv)+cwv
 =-w\Delta v-2Dw\cdot Dv+v(-\Delta w+cw)
 =-w^{-1}\operatorname{div}(w^2Dv).
\]
Thus $u$ solves the first equation exactly when $v=u/w$ solves the
second with $a=w^2$.
Conversely, given a smooth positive $a$, set $w=\sqrt a$ and
$c=\Delta w/w$. Then $-\Delta w+cw=0$, so the same identity proves
that $u=\sqrt a\,v$ solves the potential equation.
These are local identities wherever $w>0$; uniform bounds on $a$
are additional requirements for a uniform elliptic theory.
\end{proof}
\Needspace{45mm}
\PDEAnnotation{正解作为新的权重}{除以一个正解后，零阶项被吸收到权重 $w^2$ 中。关键是展开 $-\Delta(wv)$ 时保留两倍交叉项；反向构造的势是 $\Delta\sqrt a/\sqrt a$。}

\Needspace{30mm}

\Needspace{45mm}
\PDEAnnotation{负势可以由梯度能量吸收}{零阶系数略微为负并不立即破坏可解性。Poincar\'e 不等式把函数平方项转回梯度项；负势大小的阈值因此与区域几何有关。}
\paragraph{2. A small negative potential.}
For $L=-\partial_j(a^{ij}\partial_i)+c$, find $\mu>0$ so that
$c\geq-\mu$ ensures Lax--Milgram.
\begin{proof}[Solution]
Let $\|u\|_2\leq C_P\|Du\|_2$ on $H^1_0(U)$ and choose
$\mu=\theta/(2C_P^2)$. Then
\[
 B[u,u]\geq\theta\|Du\|_2^2-\mu\|u\|_2^2
 \geq\tfrac{\theta}{2}\|Du\|_2^2\geq C\|u\|_{H^1}^2.
\]
Bounded coefficients give boundedness of $B$. Apply
Theorem~\ref{thm:6-lax-milgram}. The permitted negative part depends on
both ellipticity and the domain's Poincar\'e constant.
\end{proof}

\Needspace{30mm}
\paragraph{3. The clamped biharmonic problem.}
For $f\in L^2(U)$, solve
$\Delta^2u=f$, with $u=\partial_\nu u=0$, weakly in $H^2_0(U)$.
\begin{proof}[Solution]
Use $B_2[u,v]=\int\Delta u\,\Delta v$ and $\ell(v)=\int fv$.
For $u\in C^\infty_c(U)$, integration by parts twice gives
\[
 \int(\Delta u)^2=\sum_{i,j}\int u_{ii}u_{jj}
                  =\sum_{i,j}\int u_{ij}^2.
\]
Approximation extends this to $H^2_0(U)$.
Since $u$ and each $u_i$ belong to $H^1_0(U)$, Poincar\'e applied twice
gives $\|u\|_{H^2}\leq C\|D^2u\|_2$.
Thus $B_2$ is bounded and coercive on $H^2_0$ and $\ell$ is bounded.
Lax--Milgram gives exactly one weak solution. On smooth boundaries
the closure definition of $H^2_0$ encodes both zero traces.
\end{proof}
\Needspace{45mm}
\PDEAnnotation{四阶方程用二阶能量空间}{双调和方程分部积分两次后只需二阶弱导数。$H^2_0$ 同时编码函数及一阶迹为零，不能用只有函数零迹的 $H^2\cap H^1_0$ 随意替代。}

\Needspace{30mm}
\paragraph{4. Neumann compatibility.}
For connected $U$, prove that
$\int Du\cdot Dv=\int fv$ for every $v\in H^1(U)$ is solvable exactly
when $\int_U f=0$.
\begin{proof}[Solution]
Testing with $v=1$ proves necessity. On
$V=\{u\in H^1(U):\int_Uu=0\}$, the mean-zero Poincar\'e inequality holds.
To see this, a counterexample can be normalized to have $\|u_j\|_2=1$,
mean zero and $\|Du_j\|_2\to0$.
Rellich gives a strong $L^2$ limit, whose weak gradient is zero.
It is constant by connectedness and zero by its mean, contradicting norm one.
Lax--Milgram on $V$ gives a unique mean-zero solution.
For arbitrary $v$, subtract its mean; the removed constant contributes
zero on both sides because $\int f=0$. Thus the full weak equation holds.
The difference of any two solutions has zero gradient, so all solutions
differ by a constant.
\end{proof}

\Needspace{30mm}
\paragraph{5. Robin boundary conditions.}
Formulate and solve $-\Delta u=f$, $u+\partial_\nu u=0$.
\begin{proof}[Solution]
Use $H^1(U)$ and its trace $T$, with
\[
 B_R[u,v]=\int_U Du\cdot Dv+\int_{\partial U}Tu\,Tv,\qquad
 \ell(v)=\int_Ufv.
\]
The plus boundary term follows from $\partial_\nu u=-u$ in Green's formula.
Trace boundedness makes $B_R$ bounded. Also
\[
 \|u\|_{H^1}^2\leq C\left(\|Du\|_2^2+
                            \|Tu\|_{L^2(\partial U)}^2\right).
\]
Otherwise normalize a violating sequence to norm one. Its gradients
and traces tend to zero; Rellich and weak $H^1$ convergence yield a
constant on each component, with zero trace, hence zero.
Strong $L^2$ convergence and vanishing gradients contradict norm one.
This proves coercivity; Lax--Milgram yields a unique solution for every $f$.
\end{proof}
\Needspace{45mm}
\PDEAnnotation{边界项也能控制常数}{Neumann 能量看不见常数，所以需要相容条件并允许加常数。Robin 能量多了边界平方项，常数方向也受到控制，于是无需零均值条件且解唯一。}

\Needspace{30mm}

\Needspace{45mm}
\PDEAnnotation{自然边界与本质边界}{Dirichlet 条件先限制解和测试函数所在空间，称为本质条件；Neumann 条件由分部积分后的边界项体现，称为自然条件。更换边界类型也就更换了变分空间。}
\paragraph{6. Mixed Dirichlet--Neumann data.}
Let connected $U$ have $\partial U=\Gamma_1\cup\Gamma_2$ with disjoint
closed pieces. Impose $u=0$ on $\Gamma_1$ and $\partial_\nu u=0$ on $\Gamma_2$.
\begin{proof}[Solution]
Set $V=\{u\in H^1(U):Tu=0\text{ on }\Gamma_1\}$ and require
$\int Du\cdot Dv=\int fv$ for every $v\in V$.
The Dirichlet condition is built into $V$; the Neumann condition appears
naturally after integration by parts.
If $\Gamma_1$ has positive surface measure, Poincar\'e holds on $V$:
a normalized counterexample converges to a constant by the argument
in Problem 4, and its zero trace on $\Gamma_1$ forces that constant to vanish.
Lax--Milgram gives existence and uniqueness for every $f$.
Under the stated smoothness and disjoint closed partition, a nonempty
$\Gamma_1$ is relatively open in $\partial U$ and has positive measure.
If $\Gamma_1$ is empty, Problem 4 gives the Neumann compatibility and
constant ambiguity. If $\Gamma_2$ is empty, this is the usual Dirichlet problem.
For more general partitions that meet along an interface, the weak
formulation still makes sense, but global $H^2$ regularity is a separate issue.
\end{proof}

\Needspace{30mm}
\paragraph{7. A monotone semilinear estimate.}
Let $u\in H^1(\R^n)$ have compact support and solve
$-\Delta u+c(u)=f$ weakly, where $c$ is smooth, $c(0)=0$, $c'\geq0$,
and $c(u),f\in L^2(\R^n)$. Prove $\|D^2u\|_2\leq C\|f\|_2$.
\begin{proof}[Solution]
For fixed $h$ in direction $e_k$, use
$v=-\delta_{-h}\delta_hu\in H^1(\R^n)$, an admissible compactly
supported weak test by density.
Discrete integration by parts gives
\[
 \|\delta_hDu\|_2^2+
 \int_{\R^n}\delta_h(c(u))\,\delta_hu
 =-\int_{\R^n}f\,\delta_{-h}\delta_hu.
\]
The second term is nonnegative because
$(c(s)-c(t))(s-t)\geq0$. Moreover,
$\|\delta_{-h}\delta_hu\|_2\leq\|\partial_k\delta_hu\|_2
\leq\|\delta_hDu\|_2$.
Hence $\|\delta_hDu\|_2\leq\|f\|_2$ uniformly in $h$.
The difference quotient criterion in each direction yields $u\in H^2$ and,
for the Frobenius norm, $\|D^2u\|_2\leq\sqrt n\,\|f\|_2$.
No upper bound on $c'$ is required in this monotonicity argument.
\end{proof}
\Needspace{45mm}
\PDEAnnotation{单调性替代导数上界}{非线性项的差商乘积非负，因此可以留在左边后直接舍去。这里使用的是 $c$ 的单调性，不需要把 $c'(u)$ 当作有界系数来估计。}

\Needspace{30mm}
\paragraph{8. Bernstein's gradient estimate.}
For $Lu=-a^{ij}u_{ij}=0$, show that
$v=|Du|^2+Mu^2$ is a subsolution for large $M$ and deduce a boundary estimate.
\begin{proof}[Solution]
Differentiate the equation to get
$Lu_k=(\partial_k a^{ij})u_{ij}$. Thus
\[
 L|Du|^2
 =-2a^{ij}u_{ki}u_{kj}
       +2u_k(\partial_k a^{ij})u_{ij}
 \leq-\theta|D^2u|^2+C|Du|^2 .
\]
Also $L(u^2)=-2a^{ij}u_i u_j\leq-2\theta|Du|^2$.
Choose $2M\theta\geq C$. Then $Lv\leq0$.
The weak maximum principle gives
$\sup_U|Du|^2\leq\sup_{\partial U}(|Du|^2+Mu^2)$, hence
\[
 \|Du\|_{L^\infty(U)}
 \leq C\bigl(\|Du\|_{L^\infty(\partial U)}
               +\|u\|_{L^\infty(\partial U)}\bigr).
\]
\end{proof}

\Needspace{30mm}
\paragraph{9. A boundary barrier estimate.}
For $Lu=-a^{ij}u_{ij}=f$, $u=0$ on $\partial U$, suppose
$Lw\geq1$, $w\geq0$ on $\partial U$, and $w(x_0)=0$.
\begin{proof}[Solution]
The minimum principle gives $w\geq0$ in $U$.
Let $F=\|f\|_\infty$. Since
$L(\pm u-Fw)\leq0$ and $\pm u-Fw\leq0$ on the boundary,
comparison gives $|u|\leq Fw$.
Assuming the indicated boundary derivatives exist, evaluate at
$x_0-t\nu$ and divide by $t$ to obtain
$|\partial_\nu u(x_0)|\leq F|\partial_\nu w(x_0)|$.
All tangential derivatives of the smooth zero-boundary function $u$
vanish at $x_0$, so $|Du(x_0)|=|\partial_\nu u(x_0)|$.
This is the desired estimate, with constant $F$.
\end{proof}
\Needspace{45mm}
\PDEAnnotation{屏障把函数值变成法向斜率}{先比较得到整个区域的 $|u|\leq\|f\|_\infty w$，再沿内法线取差商。边界上 $u$ 恒为零，所以梯度只剩法向分量。}

\Needspace{30mm}

\Needspace{80mm}
\PDEAnnotation{同一唯一性可有两种证明}{能量法用积分把整个梯度消掉；最大值原理通过极值位置与边界斜率制造矛盾。前者更适合弱解，后者直接揭示正性与边界几何的作用。}
\paragraph{10. Harmonic Neumann solutions are constant.}
Give energy and maximum-principle proofs on connected $U$.
\begin{proof}[Solution]
Green's identity gives
$\int_U|Du|^2=-\int_Uu\Delta u+\int_{\partial U}u\partial_\nu u=0$,
so $Du=0$ and connectedness gives constancy.
Alternatively, if $u$ were nonconstant, its maximum would be on the
boundary by the strong maximum principle and all interior values would
be strictly smaller. Hopf's lemma at a boundary maximum would imply
$\partial_\nu u>0$, contradicting the homogeneous Neumann condition.
The latter proof uses a smooth boundary, or at least the needed interior ball.
\end{proof}

\Needspace{30mm}
\paragraph{11. Convex functions of weak solutions.}
If bounded $u\in H^1(U)$ satisfies
$\int A Du\cdot Dv=0$ for $v\in H^1_0(U)$, prove that
$w=\varphi(u)$ is a weak subsolution for smooth convex $\varphi$.
\begin{proof}[Solution]
The bounded range of $u$ makes $\varphi',\varphi''$ bounded on that range,
so $w\in H^1(U)$ and $Dw=\varphi'(u)Du$.
First take bounded $0\leq v\in H^1_0(U)$.
The product $\varphi'(u)v\in H^1_0(U)$ is an admissible test and gives
\[
 \int A Dw\cdot Dv
 =-\int \varphi''(u)v\, A Du\cdot Du\leq0 .
\]
For arbitrary $0\leq v\in H^1_0(U)$, use
$v_M=\min(v,M)$ and pass in the left side by $v_M\to v$ in $H^1$.
This proves the required inequality for every nonnegative weak test.
The function $w$ need not have zero trace.
\end{proof}
\Needspace{45mm}
\PDEAnnotation{测试函数的乘积要先检查}{若 $v$ 无界，$\varphi''(u)vDu$ 未必属于 $L^2$，不能直接宣布 $\varphi'(u)v$ 合法。先截断 $v$，证明不等式后再用 $H^1$ 收敛去掉截断。}

\Needspace{30mm}
\paragraph{12. A positive supersolution implies comparison.}
Assume $h\in C^2(U)\cap C(\overline U)$,
$h>0$ on \emph{all of} $\overline U$, and $Lh\geq0$.
Prove the zero-boundary weak maximum principle without assuming $c\geq0$.
\begin{proof}[Solution]
For $Lu\leq0$ and $u\leq0$ on $\partial U$, set $z=u/h$.
The product formula gives
\[
 Lu=hMz+zLh,\qquad
 M=-A:D^2+(b-2ADh/h)\cdot D.
\]
On $\{z>0\}$, $Mz\leq0$ and $M$ has no potential.
For each $\varepsilon>0$, the closure of $\{z>\varepsilon\}$ is compact
in $U$, because $z\leq0$ on the boundary and $\min_{\overline U}h>0$.
The drift is bounded there, so the zero-potential maximum principle
gives $z\leq\varepsilon$ on this set, which must therefore be empty.
Let $\varepsilon\downarrow0$.
Positivity only in $U$ would not suffice:
$h=\sin x$ for $L=-\partial_{xx}-1$ on $(0,\pi)$ is an interior
positive solution but vanishes at the endpoints, and comparison fails.
\end{proof}

\Needspace{30mm}
\paragraph{13. Courant's max--min principle.}
For the symmetric operator of \S\ref{sec:6-5}, show
\[
 \lambda_k=\max_{\dim S=k-1}\ \min_{\substack{u\in H^1_0(U)\cap S^\perp\\\|u\|_2=1}} B[u,u],
\]
where $S\subset H^1_0(U)$ and orthogonality is in $L^2(U)$.
\begin{proof}[Solution]
For any such $S$, linear algebra supplies a nonzero vector in
$\operatorname{span}(w_1,\ldots,w_k)\cap S^\perp$.
Normalize it. Formula~\eqref{eq:6-energy-spectrum} gives its energy
at most $\lambda_k$, so the inner minimum is at most $\lambda_k$.
For $S=\operatorname{span}(w_1,\ldots,w_{k-1})$, every admissible $u$
has only modes $j\geq k$, and the same formula gives
$B[u,u]\geq\lambda_k$. Equality occurs at $w_k$.
For any fixed $S$, the minimum exists by the direct method; its
orthogonality constraints are preserved by strong $L^2$ convergence.
\end{proof}
\Needspace{45mm}
\PDEAnnotation{极小极大中的正交空间}{这里去掉的是前若干模态在 $L^2$ 中的方向。任意少于 $k$ 个线性约束都不可能消灭前 $k$ 个模态的全部组合，这就是一半不等式的有限维理由。}

\Needspace{30mm}

\Needspace{45mm}
\PDEAnnotation{非对称公式中的权重来自伴随}{不能只对 $Lu/u$ 做普通空间平均。与正的伴随特征函数配对后，正确权重是 $u\phi^*$；分部积分才把这个加权平均固定为主特征值。}
\paragraph{14. A principal-eigenvalue max--min formula.}
For the nonsymmetric operator, prove
\[
 \lambda_1=\sup_{\substack{u\in C^\infty(\overline U),\
                  u>0\text{ in }U\\u=0\text{ on }\partial U}}
                  \ \inf_{x\in U}\frac{Lu(x)}{u(x)}.
\]
\begin{proof}[Solution]
Let $\phi^*>0$ be the adjoint eigenfunction from
Lemma~\ref{lem:6-positive-adjoint}. Both functions have zero trace,
so Green's identity, or the weak adjoint identity, gives
\[
 \int_U\phi^*Lu=\int_UuL^*\phi^*
                =\lambda_1\int_Uu\phi^*.
\]
Thus the weighted average of $Lu/u$ with positive weight $u\phi^*$
is $\lambda_1$. Its infimum cannot exceed this average.
If the infimum is $-\infty$ the inequality is immediate; otherwise
integrating its pointwise lower bound proves it directly.
Taking $u=\phi$ achieves equality everywhere.
A constant shift handles potentials of arbitrary bounded sign.
\end{proof}

\Needspace{30mm}
\paragraph{15. Hadamard's domain-variation formula.}
Let $-\Delta w=\lambda w$ on a smoothly moving domain $U(\tau)$,
$w=0$ on its boundary and $\int_{U(\tau)}w^2=1$.
Assume a smooth eigenvalue/eigenfunction branch, and denote boundary
velocity by $\mathbf v$.
\begin{proof}[Solution]
Use the Eulerian derivative $\dot w$ at a fixed spatial point.
Differentiating $w(x(\tau),\tau)=0$ on the moving boundary gives
\[
 \dot w=-Dw\cdot\mathbf v
       =-\partial_\nu w\,(\mathbf v\cdot\nu).
\]
Reynolds' formula for normalization yields
$2\int_Uw\dot w+\int_{\partial U}w^2\mathbf v\cdot\nu=0$,
hence $\int_Uw\dot w=0$.
Since $\lambda=\int_U|Dw|^2$, differentiate once more:
\begin{align*}
 \dot\lambda
 &=2\int_U Dw\cdot D\dot w+
             \int_{\partial U}|Dw|^2\mathbf v\cdot\nu\\
 &=2\lambda\int_Uw\dot w+
      2\int_{\partial U}\partial_\nu w\,\dot w+
             \int_{\partial U}(\partial_\nu w)^2\mathbf v\cdot\nu\\
 &=-\int_{\partial U}(\partial_\nu w)^2\mathbf v\cdot\nu.
\end{align*}
Tangential derivatives vanish because $w$ is zero on the boundary.
The smooth-branch hypothesis is essential; a multiple eigenvalue need
not admit a single arbitrarily chosen differentiable eigenfunction branch.
\end{proof}
\Needspace{45mm}
\PDEAnnotation{扩张区域使特征值下降}{若边界向外移动，$\mathbf v\cdot\nu>0$，公式给出 $\dot\lambda<0$。负号来自移动边界条件 $\dot w=-\partial_\nu w\,\mathbf v\cdot\nu$，不是凭直觉添加。}

\Needspace{30mm}

\Needspace{45mm}
\PDEAnnotation{辐射条件选择向外传播}{时间因子取 $e^{-i\sigma t}$ 时，$e^{i\sigma r}/r$ 描述向外传播的球面波。条件中的减号与这一时间约定配套；平面波不会在所有方向都向外辐射。}
\paragraph{16. Helmholtz waves and radiation.}
Let real $\sigma\ne0$ and $\lambda=\sigma^2$.
\begin{proof}[Solution]
(a) For $|\omega|=1$,
$-\Delta e^{i\sigma\omega\cdot x}=\sigma^2e^{i\sigma\omega\cdot x}$.
Thus $e^{i\sigma(\omega\cdot x-t)}$ solves the wave equation.

(b) In dimension three put
$\Phi(x)=e^{i\sigma r}/(4\pi r)$, $r=|x|$.
The radial identity $\Delta g=g''+2g'/r$ shows
$(-\Delta-\sigma^2)\Phi=0$ for $r>0$.
At $r=\varepsilon$,
\[
 -\int_{\partial B_\varepsilon}\partial_r\Phi\,dS
 =e^{i\sigma\varepsilon}(1-i\sigma\varepsilon)\longrightarrow1.
\]
Integration by parts against a smooth test function outside
$B_\varepsilon$ therefore gives
$(-\Delta-\sigma^2)\Phi=\delta_0$; the other shrinking-sphere term and
the volume integral of $\sigma^2\Phi$ tend to zero.

(c) For $\Phi$,
$r(\Phi_r-i\sigma\Phi)=-e^{i\sigma r}/(4\pi r)\to0$ uniformly in angle.
For the plane wave at $x=r\vartheta$,
\[
 r(w_r-i\sigma w)
 =i\sigma r(\omega\cdot\vartheta-1)e^{i\sigma r\omega\cdot\vartheta},
\]
whose modulus diverges on every ray with $\omega\cdot\vartheta\ne1$.
Thus the plane wave does not satisfy the Sommerfeld condition.
\end{proof}

\Needspace{30mm}
\paragraph{17. Uniqueness of an entire radiating solution.}
Suppose $w$ solves $-\Delta w=\sigma^2w$ throughout $\R^3$ and satisfies
$r(w_r-i\sigma w)\to0$ \emph{uniformly in angle}. Prove $w=0$.
\begin{proof}[Solution]
Green's identity and the real coefficient $\sigma^2$ imply
\[
 0=\int_{\partial B_R}(\overline w\,w_r-w\,\overline{w_r})\,dS.
\]
The imaginary flux is zero, so expansion of the square gives
\[
 \int_{\partial B_R}(|w_r|^2+\sigma^2|w|^2)\,dS
 =\int_{\partial B_R}|w_r-i\sigma w|^2\,dS=o(1).
\]
Uniform radiation justifies the last step, since the sphere has area
$4\pi R^2$ and the integrand is $o(R^{-2})$ uniformly.
For a fixed $x_0$ and $R>|x_0|$, Green's representation with the
outgoing kernel $\Phi(x-x_0)$ gives
\[
 w(x_0)=\int_{\partial B_R}
   \bigl(\Phi\,w_r-w\,\partial_r\Phi\bigr)\,dS .
\]
Here $\partial_r$ on the kernel means the outward derivative on
$\partial B_R$, not radial differentiation about $x_0$.
For large $R$, both $\Phi(x-x_0)$ and this derivative are $O(R^{-1})$,
so their $L^2(\partial B_R)$ norms remain bounded.
Cauchy--Schwarz and the vanishing flux norm above imply $w(x_0)=0$.
Since $x_0$ was arbitrary, $w\equiv0$.
\end{proof}
\Needspace{45mm}
\PDEAnnotation{辐射条件还包含一致性}{从逐方向极限跳到球面积分极限需要控制角变量。这里明确采用对方向一致的 Sommerfeld 条件；基本解在原点有奇性，所以并不违背“全空间光滑辐射解只能为零”。}
